\chapter{Business Requirements Specification}
\label{chap:business-requirements}

\section*{Integration with existing chapters}
\addcontentsline{toc}{section}{Integration with existing chapters}

This chapter provides the business governance layer around the technical
specifications in earlier chapters:

\begin{itemize}
  \item \textbf{Chapter~\ref{sec:sim-roadmap}} (Project Roadmap): BR-1.x strategic 
    requirements align with the MoSCoW classification in \cref{tab:sim-moscow}. 
    High-priority BRs map to Must milestones (M01--M07); Medium-priority BRs map 
    to Should (S01--S02); Low-priority BRs map to Could (C01).
  \item \textbf{Chapter~\ref{ch:agent-impl}} (Agent Instructions): BR-2.x generation 
    requirements provide acceptance criteria for Tasks S1--S4 (schema extraction, 
    taxonomy building, generation, evaluation).
  \item \textbf{Chapter~\ref{chap:enterprise-enhancements}} (Enterprise Enhancements): 
    BR-3.x quality requirements are implemented by Enhancement~3 (Closed-Loop 
    Validation); BR-4.x governance requirements by Enhancement~4 (Workflow Engine).
  \item \textbf{Table~\ref{tab:simula-fmea}} (FMEA): BR-3.3 (Privacy) and BR-3.4 
    (Bias) address the ``Critic Collusion'' and ``Schema Contraction'' failure modes.
\end{itemize}

\noindent
The traceability matrix at the end of this chapter maps each BR to the enhancement 
and regulatory framework that governs it.

\section*{Scope of this chapter}
\addcontentsline{toc}{section}{Scope of this chapter}

This chapter defines the complete business requirements for enterprise deployment
of the Simula synthetic data framework. Requirements are organized across five
categories: strategic preparation, generation phase, quality assurance, governance
and compliance, and integration infrastructure. Each requirement includes regulatory
citations where applicable.

\section{Requirements Overview}

\subsection{Requirement Identification Convention}

Requirements follow the identification pattern \texttt{BR-<category>.<sequence>}:

\begin{table}[htbp]
\centering
\caption{Business requirement categories.}
\label{tab:br-categories}
\footnotesize
\begin{tabularx}{\textwidth}{@{}l l X l@{}}
\toprule
\textbf{Category} & \textbf{Prefix} & \textbf{Scope} & \textbf{Count} \\
\midrule
Strategic \& Preparation & BR-1.x & Use case definition, governance integration & 5 \\
Generation Phase & BR-2.x & Taxonomy, configuration, metadata & 5 \\
Quality Assurance & BR-3.x & Validation, privacy, bias detection & 5 \\
Governance \& Compliance & BR-4.x & Access control, audit, transparency & 5 \\
Integration \& Infrastructure & BR-5.x & ML pipelines, storage, deployment & 5 \\
\bottomrule
\end{tabularx}
\end{table}

\subsection{Priority Classification}

Requirements are classified using MoSCoW prioritization:
\begin{itemize}
  \item \textbf{High}: Must have for initial deployment
  \item \textbf{Medium}: Should have for production maturity
  \item \textbf{Low}: Could have for optimization
\end{itemize}


\section{BR-1: Strategic and Preparation Phase}

\subsection{BR-1.1: Use Case Definition}

\begin{tcolorbox}[colback=blue!5,colframe=blue!70,title=BR-1.1: Clear Use Case Definition]
\textbf{Priority}: High

\textbf{Description}: Organizations must clearly define synthetic data use cases
and objectives before initiating generation projects. The synthetic data market
is projected to reach \$3.4 billion by 2031\footnote{Market Research Future,
``Synthetic Data Generation Market Report,'' 2024}, but success depends on
starting with high-value, well-scoped use cases.

\textbf{Typical Use Cases}:
\begin{itemize}
  \item Privacy-sensitive domain data substitution
  \item Long-tail scenario coverage augmentation
  \item Adversarial testing data generation
  \item Training data augmentation for rare events
  \item Model validation in regulated environments
\end{itemize}

\textbf{Acceptance Criteria}:
\begin{enumerate}
  \item Documented use case with clear business objective
  \item Identified target domain and data modality
  \item Defined success metrics and measurement approach
  \item Stakeholder sign-off on scope and expectations
\end{enumerate}

\textbf{SAP Integration}: Align use cases with SAP AI Core deployment scenarios
(Text-to-SQL, document extraction, financial analytics).
\end{tcolorbox}

\subsection{BR-1.2: Value-Risk Assessment}

\begin{tcolorbox}[colback=blue!5,colframe=blue!70,title=BR-1.2: Structured Value-Risk Assessment]
\textbf{Priority}: High

\textbf{Description}: Before initiating synthetic data projects, organizations
must conduct structured value-risk assessments. The UK FCA Synthetic Data Expert
Group explicitly identifies clear objective definition and value-risk analysis
as foundational to project success\footnote{FCA, ``Synthetic Data Expert Group
Report,'' Section 3.2, 2024}.

\textbf{Assessment Dimensions}:
\begin{itemize}
  \item \textbf{Value}: Cost savings, time-to-market, compliance enablement
  \item \textbf{Risk}: Privacy leakage, bias amplification, quality degradation
  \item \textbf{Feasibility}: Technical capability, resource availability
  \item \textbf{Regulatory}: Jurisdiction-specific compliance requirements
\end{itemize}

\textbf{Acceptance Criteria}:
\begin{enumerate}
  \item Completed value-risk assessment matrix
  \item Identified risk mitigation strategies
  \item Documented go/no-go decision criteria
  \item Executive sponsor approval
\end{enumerate}
\end{tcolorbox}

\subsection{BR-1.3: Governance Framework Integration}

\begin{tcolorbox}[colback=blue!5,colframe=blue!70,title=BR-1.3: Existing Governance Framework Integration]
\textbf{Priority}: High

\textbf{Description}: In the absence of dedicated synthetic data governance
frameworks, organizations should integrate synthetic data management into existing
model risk management, data governance, and AI ethics frameworks rather than
building from scratch.

\textbf{Integration Points}:
\begin{itemize}
  \item Model Risk Management (MRM) --- treat synthetic data as model input
  \item Data Governance --- apply data quality and lineage standards
  \item AI Ethics Framework --- extend bias and fairness requirements
  \item Information Security --- apply data classification policies
\end{itemize}

\textbf{Acceptance Criteria}:
\begin{enumerate}
  \item Synthetic data policies integrated into existing governance documentation
  \item Roles and responsibilities defined within existing RACI matrices
  \item Approval workflows configured in existing governance tools
  \item Training materials updated for affected stakeholders
\end{enumerate}

\textbf{SAP Integration}: Leverage SAP GRC (Governance, Risk, Compliance) for
policy management and workflow orchestration.
\end{tcolorbox}

\subsection{BR-1.4: Cross-Functional Team Formation}

\begin{tcolorbox}[colback=blue!5,colframe=blue!70,title=BR-1.4: Cross-Functional Team Organization]
\textbf{Priority}: Medium

\textbf{Description}: Synthetic data generation and usage requires multi-stakeholder
collaboration. The World Economic Forum notes that synthetic data governance is
typically constructed by multiple distinct teams working together\footnote{WEF,
``Synthetic Data for Financial Services,'' p.~18, 2024}.

\textbf{Required Roles}:
\begin{itemize}
  \item \textbf{Data Science Team}: Technical implementation, model training
  \item \textbf{Domain Experts}: Taxonomy validation, complexity calibration
  \item \textbf{Legal/Compliance}: Privacy risk assessment, regulatory alignment
  \item \textbf{IT/Platform}: Infrastructure provisioning, security controls
  \item \textbf{Business Stakeholders}: Use case definition, acceptance criteria
\end{itemize}

\textbf{Acceptance Criteria}:
\begin{enumerate}
  \item Documented team structure with named individuals
  \item Defined communication and escalation paths
  \item Regular sync cadence established
  \item Clear decision-making authority matrix
\end{enumerate}
\end{tcolorbox}

\subsection{BR-1.5: ROI Baseline Establishment}

\begin{tcolorbox}[colback=blue!5,colframe=blue!70,title=BR-1.5: Investment Return Baseline]
\textbf{Priority}: Medium

\textbf{Description}: Organizations should establish clear ROI measurement
baselines. Industry data indicates typical benefits include: AI model development
speed improvements up to 40\%, compliance and data processing cost reductions
of 50--60\%\footnote{Gartner, ``Synthetic Data Adoption Trends,'' 2025}.

\textbf{Measurement Dimensions}:
\begin{itemize}
  \item Time-to-data: Days from request to available training data
  \item Data acquisition cost: Cost per training example
  \item Model performance: Accuracy, F1, or domain-specific metrics
  \item Compliance cost: Audit preparation, privacy review cycles
\end{itemize}

\textbf{Acceptance Criteria}:
\begin{enumerate}
  \item Baseline metrics captured before synthetic data adoption
  \item Target improvement thresholds defined
  \item Measurement methodology documented
  \item Reporting cadence established
\end{enumerate}
\end{tcolorbox}


\section{BR-2: Generation Phase}

\subsection{BR-2.1: Domain Taxonomy Construction and Validation}

\begin{tcolorbox}[colback=green!5,colframe=green!70,title=BR-2.1: Taxonomy Construction and Validation]
\textbf{Priority}: High

\textbf{Description}: Simula's global diversification step requires constructing
deep hierarchical taxonomies for target domains. Business requirements include
domain expert participation in taxonomy review and version management as domain
knowledge evolves.

\textbf{Requirements}:
\begin{itemize}
  \item Expert review interface for taxonomy inspection and correction
  \item Version control for taxonomy artifacts
  \item Coverage gap identification and reporting
  \item Domain vocabulary injection mechanism
\end{itemize}

\textbf{Acceptance Criteria}:
\begin{enumerate}
  \item Taxonomy review workflow with expert approval gates
  \item Taxonomy versioning with diff visualization
  \item Coverage reports showing sampling breadth and depth
  \item Domain-specific node validation rules
\end{enumerate}

\textbf{SAP Integration}: Store taxonomy versions in SAP HANA Cloud with metadata
tracking. Integrate with SAP Signavio for process-aware taxonomy construction.
\end{tcolorbox}

\subsection{BR-2.2: Controllable Parameter Configuration}

\begin{tcolorbox}[colback=green!5,colframe=green!70,title=BR-2.2: Parameter Configuration Interface]
\textbf{Priority}: High

\textbf{Description}: Provide business user-friendly interfaces for independent
control of generation parameters without requiring technical expertise.

\textbf{Controllable Parameters}:
\begin{itemize}
  \item \textbf{Global Coverage}: Taxonomy sampling breadth and depth
  \item \textbf{Local Diversity}: Variants per concept node
  \item \textbf{Complexity Distribution}: Easy/medium/hard ratio curve
  \item \textbf{Quality Thresholds}: Critic pass rate, diversity minimums
\end{itemize}

\textbf{Acceptance Criteria}:
\begin{enumerate}
  \item Visual configuration interface with parameter sliders
  \item Preview mode showing expected distribution
  \item Parameter validation with constraint warnings
  \item Configuration templates for common scenarios
\end{enumerate}

\textbf{SAP Integration}: Expose configuration via SAP Build Apps for low-code
workflow construction. Store configurations in SAP HANA Cloud.
\end{tcolorbox}

\subsection{BR-2.3: Cost and Resource Budget Control}

\begin{tcolorbox}[colback=green!5,colframe=green!70,title=BR-2.3: Budget Control Mechanism]
\textbf{Priority}: Medium

\textbf{Description}: Synthetic data generation involves inference model API
costs. Organizations require budget control mechanisms to manage expenses.

\textbf{Requirements}:
\begin{itemize}
  \item Per-project and per-department generation quotas
  \item Single-task cost ceiling configuration
  \item Resource consumption estimation before execution
  \item Low-cost ``draft mode'' for taxonomy design validation
\end{itemize}

\textbf{Acceptance Criteria}:
\begin{enumerate}
  \item Budget allocation interface by cost center
  \item Real-time cost tracking during generation
  \item Alerts when approaching budget limits
  \item Cost-per-example reporting
\end{enumerate}

\textbf{SAP Integration}: Integrate with SAP AI Core resource quotas and
SAP Analytics Cloud for cost visualization.
\end{tcolorbox}

\subsection{BR-2.4: Metadata and Provenance Recording}

\begin{tcolorbox}[colback=green!5,colframe=green!70,title=BR-2.4: Complete Metadata Capture]
\textbf{Priority}: High

\textbf{Description}: Each generation task must automatically record complete
metadata for audit compliance and reproducibility.

\textbf{Required Metadata}:
\begin{itemize}
  \item Reasoning model version and configuration
  \item Taxonomy structure snapshot
  \item Sampling strategy parameters
  \item Generation timestamp and executor identity
  \item Random seeds for reproducibility
\end{itemize}

\textbf{Acceptance Criteria}:
\begin{enumerate}
  \item Automatic metadata capture on every generation run
  \item Metadata stored with tamper-evident integrity
  \item Query interface for metadata retrieval
  \item Export capability for audit packages
\end{enumerate}

\textbf{SAP Integration}: Store provenance in SAP HANA Cloud with blockchain-style
hash chaining. Integrate with SAP Data Intelligence for lineage visualization.
\end{tcolorbox}

\subsection{BR-2.5: Incremental Generation and Dataset Evolution}

\begin{tcolorbox}[colback=green!5,colframe=green!70,title=BR-2.5: Incremental Generation Support]
\textbf{Priority}: Low

\textbf{Description}: Support incremental generation on existing datasets when
taxonomies evolve, avoiding redundant regeneration.

\textbf{Requirements}:
\begin{itemize}
  \item Detect new/modified taxonomy nodes
  \item Generate only for changed coverage regions
  \item Merge incremental examples with existing dataset
  \item Version-to-version diff analysis
\end{itemize}

\textbf{Acceptance Criteria}:
\begin{enumerate}
  \item Taxonomy diff computation between versions
  \item Targeted generation for changed nodes only
  \item Dataset merge with deduplication
  \item Generation efficiency metrics (examples saved)
\end{enumerate}
\end{tcolorbox}


\section{BR-3: Quality Assurance and Validation}

\subsection{BR-3.1: Intrinsic Quality Assessment}

\begin{tcolorbox}[colback=yellow!5,colframe=yellow!70,title=BR-3.1: Standardized Quality Reporting]
\textbf{Priority}: High

\textbf{Description}: Simula's built-in quality assessment capabilities must be
productized into standardized reports for business consumption.

\textbf{Required Reports}:
\begin{itemize}
  \item \textbf{Coverage Report}: Taxonomy branches sampled vs.\ omitted
  \item \textbf{Complexity Distribution}: Generated data complexity vs.\ target
  \item \textbf{Diversity Metrics}: Intra-concept sample similarity analysis
  \item \textbf{Critic Pass Rates}: Acceptance rates by taxonomy region
\end{itemize}

\textbf{Acceptance Criteria}:
\begin{enumerate}
  \item Automated report generation on every dataset
  \item Visual dashboards with drill-down capability
  \item Threshold alerts for quality degradation
  \item Historical trend analysis
\end{enumerate}

\textbf{SAP Integration}: Publish reports to SAP Analytics Cloud. Configure
alerts via SAP Alert Notification Service.
\end{tcolorbox}

\subsection{BR-3.2: Downstream Task Validation}

\begin{tcolorbox}[colback=yellow!5,colframe=yellow!70,title=BR-3.2: Downstream Model Validation]
\textbf{Priority}: High

\textbf{Description}: Generated datasets must be validated on downstream models
to verify practical utility. Recommended approach: ``Train-Synthetic-Test-Real''
(TSTR) methodology.

\textbf{Validation Protocol}:
\begin{enumerate}
  \item Train model on synthetic data
  \item Evaluate on held-out real data
  \item Compute synthetic-to-real performance gap
  \item Identify systematic failure patterns
\end{enumerate}

\textbf{Acceptance Criteria}:
\begin{enumerate}
  \item Automated TSTR pipeline integration
  \item Performance gap quantification by taxonomy region
  \item Failure pattern analysis reports
  \item Threshold-based quality gates
\end{enumerate}

\textbf{SAP Integration}: Execute validation via SAP AI Core training pipelines.
Store results in SAP HANA Cloud for trend analysis.
\end{tcolorbox}

\subsection{BR-3.3: Privacy Risk Assessment}

\begin{tcolorbox}[colback=yellow!5,colframe=yellow!70,title=BR-3.3: Privacy Risk Evaluation]
\textbf{Priority}: High

\textbf{Description}: Even in Simula's seedless mode, generated data may exhibit
unintended similarity to real data seen during reasoning model training. Privacy
risk assessment mechanisms are essential. IEEE emphasizes that privacy risk
metrics and mitigation techniques are core governance components\footnote{IEEE,
``Synthetic Data: Generation, Validation, and Governance,'' Section 4.3, 2025}.

\textbf{Assessment Mechanisms}:
\begin{itemize}
  \item Membership inference attack testing
  \item Similarity scanning against known real datasets
  \item Differential privacy budget tracking
  \item Re-identification risk scoring
\end{itemize}

\textbf{Acceptance Criteria}:
\begin{enumerate}
  \item Automated privacy risk scoring on generated datasets
  \item Configurable risk thresholds with blocking gates
  \item Audit-ready privacy assessment reports
  \item Remediation workflow for high-risk samples
\end{enumerate}

\textbf{SAP Integration}: Integrate with SAP Data Privacy Integration service
for sensitive data detection.
\end{tcolorbox}

\subsection{BR-3.4: Bias Detection and Mitigation}

\begin{tcolorbox}[colback=yellow!5,colframe=yellow!70,title=BR-3.4: Systematic Bias Detection]
\textbf{Priority}: Medium

\textbf{Description}: Establish systematic bias detection processes to identify
and mitigate inherited biases from reasoning models.

\textbf{Detection Approaches}:
\begin{itemize}
  \item Distribution analysis on sensitive attributes
  \item Reasoning model bias inheritance detection
  \item Over/under-sampling recommendations
  \item Fairness metric computation
\end{itemize}

\textbf{Acceptance Criteria}:
\begin{enumerate}
  \item Automated bias scanning on generation
  \item Fairness metric dashboards
  \item Bias mitigation tool integration
  \item Bias audit trail for compliance
\end{enumerate}
\end{tcolorbox}

\subsection{BR-3.5: Continuous Monitoring and Drift Detection}

\begin{tcolorbox}[colback=yellow!5,colframe=yellow!70,title=BR-3.5: Quality Drift Monitoring]
\textbf{Priority}: Medium

\textbf{Description}: Establish continuous monitoring to detect quality changes
due to reasoning model upgrades or taxonomy drift. FCA includes ``continuous
monitoring and improvement'' as one of nine governance principles\footnote{FCA,
``Synthetic Data Expert Group Report,'' Principle 9, 2024}.

\textbf{Monitoring Capabilities}:
\begin{itemize}
  \item Periodic quality sampling and assessment
  \item Version-to-version quality drift detection
  \item Model upgrade impact analysis
  \item Automated regeneration triggers
\end{itemize}

\textbf{Acceptance Criteria}:
\begin{enumerate}
  \item Scheduled quality monitoring jobs
  \item Drift detection with statistical significance testing
  \item Alert configuration for threshold breaches
  \item Automated remediation workflow triggers
\end{enumerate}

\textbf{SAP Integration}: Configure monitoring via SAP AI Core MLOps capabilities.
Route alerts through SAP Alert Notification Service.
\end{tcolorbox}


\section{BR-4: Governance and Compliance}

\subsection{BR-4.1: Access Control and Permission Management}

\begin{tcolorbox}[colback=red!5,colframe=red!70,title=BR-4.1: Role-Based Access Control]
\textbf{Priority}: High

\textbf{Description}: Although synthetic data does not directly contain real
personal information, strict access controls remain necessary.

\textbf{Control Requirements}:
\begin{itemize}
  \item Role-based access (generator, reviewer, consumer separation)
  \item Dataset access approval workflows
  \item Additional restrictions for sensitive domains (healthcare, finance)
  \item Access logging and audit trails
\end{itemize}

\textbf{Acceptance Criteria}:
\begin{enumerate}
  \item RBAC configuration in identity provider
  \item Approval workflow for dataset access requests
  \item Domain-specific access policies
  \item Comprehensive access audit logs
\end{enumerate}

\textbf{SAP Integration}: Leverage SAP Cloud Identity Services for authentication.
Configure authorization via SAP Authorization and Trust Management Service (XSUAA).
\end{tcolorbox}

\subsection{BR-4.2: Regulatory Framework Alignment}

\begin{tcolorbox}[colback=red!5,colframe=red!70,title=BR-4.2: Compliance Framework Alignment]
\textbf{Priority}: High

\textbf{Description}: Synthetic data generation and usage must align with
applicable regulations. FCA emphasizes that generating and using synthetic data
requires consideration of data protection, non-discrimination, and bias
regulatory obligations\footnote{FCA, ``Synthetic Data Expert Group Report,''
Section 5.1, 2024}.

\textbf{Regulatory Considerations}:
\begin{itemize}
  \item \textbf{GDPR}: Demonstrate anonymization status of synthetic data
  \item \textbf{HIPAA}: Comply with protected health information rules
  \item \textbf{Industry-Specific}: Financial services (FCA), healthcare (FDA)
  \item \textbf{AI Regulations}: EU AI Act risk categorization
\end{itemize}

\textbf{Acceptance Criteria}:
\begin{enumerate}
  \item Documented compliance position for each applicable regulation
  \item Legal review sign-off on synthetic data usage
  \item Regulatory change monitoring process
  \item Jurisdiction-specific policy variations
\end{enumerate}

\textbf{SAP Integration}: Integrate with SAP Compliance and GRC solutions for
regulatory tracking.
\end{tcolorbox}

\subsection{BR-4.3: Audit Logging and Traceability}

\begin{tcolorbox}[colback=red!5,colframe=red!70,title=BR-4.3: Complete Audit Trail]
\textbf{Priority}: High

\textbf{Description}: Establish complete audit tracking. WEF emphasizes that
enterprises must prioritize oversight and compliance, building robust traceability
and provenance systems\footnote{WEF, ``Synthetic Data for Financial Services,''
p.~24, 2024}.

\textbf{Audit Requirements}:
\begin{itemize}
  \item ``Who, when, why, what'' for every data access
  \item Complete usage chain (generation $\to$ training $\to$ deployment)
  \item Compliance review decision records
  \item Long-term audit log retention
\end{itemize}

\textbf{Acceptance Criteria}:
\begin{enumerate}
  \item Tamper-evident audit log storage
  \item Query interface for audit investigations
  \item Retention policy enforcement
  \item Audit export for regulatory submissions
\end{enumerate}

\textbf{SAP Integration}: Utilize SAP Audit Log Service for tamper-evident logging.
Store long-term archives in SAP Object Store with retention policies.
\end{tcolorbox}

\subsection{BR-4.4: Transparency and Explainability}

\begin{tcolorbox}[colback=red!5,colframe=red!70,title=BR-4.4: External Transparency]
\textbf{Priority}: Medium

\textbf{Description}: Organizations must be able to explain synthetic data usage
to external stakeholders (customers, regulators, partners). This addresses FCA
principles of ``transparency'' and ``explainability''\footnote{FCA, ``Synthetic
Data Expert Group Report,'' Principles 1--2, 2024}.

\textbf{Transparency Requirements}:
\begin{itemize}
  \item Documentation of synthetic nature and generation method
  \item Explanation of synthetic vs.\ real data differences
  \item Justification for synthetic data usage
  \item Model card integration for downstream models
\end{itemize}

\textbf{Acceptance Criteria}:
\begin{enumerate}
  \item Synthetic data documentation templates
  \item Stakeholder communication guidelines
  \item Model card integration with synthetic data lineage
  \item FAQ materials for common inquiries
\end{enumerate}
\end{tcolorbox}

\subsection{BR-4.5: Incident Response and Remediation}

\begin{tcolorbox}[colback=red!5,colframe=red!70,title=BR-4.5: Incident Response Process]
\textbf{Priority}: Medium

\textbf{Description}: Establish incident response processes for synthetic data
related security events.

\textbf{Incident Scenarios}:
\begin{itemize}
  \item Generated data inadvertently contains real information
  \item Systematic bias causes discriminatory model outputs
  \item Generated data is misused or leaked
  \item Quality degradation impacts downstream systems
\end{itemize}

\textbf{Acceptance Criteria}:
\begin{enumerate}
  \item Incident classification and severity matrix
  \item Responsible party identification
  \item Escalation and notification paths
  \item Remediation playbooks
\end{enumerate}

\textbf{SAP Integration}: Integrate with existing SAP incident management workflows.
Configure alerts through SAP Alert Notification Service.
\end{tcolorbox}


\section{BR-5: Integration and Infrastructure}

\subsection{BR-5.1: ML Pipeline Integration}

\begin{tcolorbox}[colback=purple!5,colframe=purple!70,title=BR-5.1: MLOps Pipeline Integration]
\textbf{Priority}: High

\textbf{Description}: Simula must integrate as an orchestratable component within
existing ML pipelines.

\textbf{Integration Requirements}:
\begin{itemize}
  \item API-driven generation task triggering
  \item Integration with major MLOps platforms
  \item Automatic dataset push to training storage
  \item Pipeline status and progress reporting
\end{itemize}

\textbf{Acceptance Criteria}:
\begin{enumerate}
  \item REST/gRPC API for generation tasks
  \item Pre-built integrations (Kubeflow, MLflow, Vertex AI)
  \item Webhook notifications for completion events
  \item Pipeline artifact registration
\end{enumerate}

\textbf{SAP Integration}: Primary integration with SAP AI Core. Secondary
integrations with SAP Data Intelligence and SAP BTP Kyma Runtime for
containerized execution.
\end{tcolorbox}

\subsection{BR-5.2: Data Storage and Version Management}

\begin{tcolorbox}[colback=purple!5,colframe=purple!70,title=BR-5.2: Unified Data Governance]
\textbf{Priority}: Medium

\textbf{Description}: Synthetic datasets should follow the same governance
policies as real datasets within organizational data infrastructure.

\textbf{Storage Requirements}:
\begin{itemize}
  \item Version-controlled data lake/warehouse storage
  \item Dataset labeling, search, and reuse support
  \item Data catalog integration
  \item Retention policy enforcement
\end{itemize}

\textbf{Acceptance Criteria}:
\begin{enumerate}
  \item Storage in enterprise data lake with versioning
  \item Metadata registration in data catalog
  \item Search and discovery interface
  \item Automated retention enforcement
\end{enumerate}

\textbf{SAP Integration}: Store datasets in SAP HANA Cloud or SAP Datasphere.
Register metadata in SAP Data Intelligence catalog.
\end{tcolorbox}

\subsection{BR-5.3: Elastic Compute Resource Scheduling}

\begin{tcolorbox}[colback=purple!5,colframe=purple!70,title=BR-5.3: Compute Resource Management]
\textbf{Priority}: Medium

\textbf{Description}: Synthetic data generation is compute-intensive. Support
dynamic resource allocation based on workload.

\textbf{Scheduling Requirements}:
\begin{itemize}
  \item Dynamic resource allocation by task size
  \item Batch and real-time generation modes
  \item Priority queue management
  \item Cost-optimized resource selection
\end{itemize}

\textbf{Acceptance Criteria}:
\begin{enumerate}
  \item Auto-scaling based on queue depth
  \item Priority scheduling with preemption
  \item Resource utilization dashboards
  \item Cost optimization recommendations
\end{enumerate}

\textbf{SAP Integration}: Leverage SAP AI Core auto-scaling. Use SAP BTP
Kubernetes runtime for containerized workloads.
\end{tcolorbox}

\subsection{BR-5.4: Multi-Environment Deployment Support}

\begin{tcolorbox}[colback=purple!5,colframe=purple!70,title=BR-5.4: Deployment Flexibility]
\textbf{Priority}: Low

\textbf{Description}: Support multiple deployment models to accommodate varying
security requirements across organizations.

\textbf{Deployment Models}:
\begin{itemize}
  \item \textbf{Public Cloud SaaS}: Standard scenarios
  \item \textbf{Private Deployment}: High-security requirements
  \item \textbf{VPC Deployment}: Regulated industry compliance
  \item \textbf{Hybrid}: Combination of above
\end{itemize}

\textbf{Acceptance Criteria}:
\begin{enumerate}
  \item Documented deployment architectures for each model
  \item Security certifications per deployment type
  \item Network isolation configuration guides
  \item Data residency compliance options
\end{enumerate}

\textbf{SAP Integration}: Leverage SAP BTP multi-environment support. Configure
SAP Private Link for secure connectivity.
\end{tcolorbox}


\section{Requirements Traceability Matrix}

\begin{table}[htbp]
\centering
\caption{Requirements to enhancement traceability.}
\label{tab:req-traceability}
\scriptsize
\begin{tabularx}{\textwidth}{@{}l l l X@{}}
\toprule
\textbf{Requirement} & \textbf{Priority} & \textbf{Enhancement} & \textbf{Regulatory Reference} \\
\midrule
BR-1.1 & High & --- & --- \\
BR-1.2 & High & --- & FCA Principle 3 \\
BR-1.3 & High & E4 & WEF Governance \\
BR-1.4 & Medium & --- & WEF Cross-functional \\
BR-1.5 & Medium & E3 & --- \\
\midrule
BR-2.1 & High & E2 & --- \\
BR-2.2 & High & E4 & --- \\
BR-2.3 & Medium & E4 & --- \\
BR-2.4 & High & E4 & FCA Principle 7 \\
BR-2.5 & Low & E4 & --- \\
\midrule
BR-3.1 & High & E3 & --- \\
BR-3.2 & High & E3 & IEEE Fidelity \\
BR-3.3 & High & E3 & IEEE Privacy, FCA Principle 6 \\
BR-3.4 & Medium & E3 & FCA Principle 8 \\
BR-3.5 & Medium & E3 & FCA Principle 9 \\
\midrule
BR-4.1 & High & E4 & FCA Principle 6 \\
BR-4.2 & High & --- & GDPR, HIPAA, FCA \\
BR-4.3 & High & E4 & WEF Traceability \\
BR-4.4 & Medium & E4 & FCA Principles 1--2 \\
BR-4.5 & Medium & --- & --- \\
\midrule
BR-5.1 & High & E3, E4 & --- \\
BR-5.2 & Medium & E4 & IEEE Provenance \\
BR-5.3 & Medium & --- & --- \\
BR-5.4 & Low & --- & --- \\
\bottomrule
\end{tabularx}
\end{table}

Legend: E2 = Domain Calibration, E3 = Closed-Loop Validation, E4 = Workflow Engine